% ============================================================================
% 10_theory.tex  --  Theory / derivation  (~1 - 1.5 pages)
%
% GOAL: The physics or derivation your part owns. Mine the matching notes/*.tex
%       listed for your part in ../../reports/README.md -- adapt and cite them,
%       do NOT re-derive from scratch. Equations use \numberwithin so they read
%       as (sec.eq).
% ============================================================================

\section{Noise model and diagnostic}
\label{sec:theory}

\subsection{Gate-level quantum channels}

\ai{An ideal gate $g$ acts on a density matrix as
$\mathcal U_g(\rho)=U_g\rho U_g^\dagger$.  In the circuit-level model the ideal
gate is followed immediately by a completely positive trace-preserving channel,}
\begin{equation}
  \widetilde{\mathcal U}_g=\mathcal E_g\circ\mathcal U_g,
  \qquad
  \mathcal E_g(\rho)=\sum_k K_k\rho K_k^\dagger,
  \qquad
  \sum_k K_k^\dagger K_k=I .
  \label{eq:kraus}
\end{equation}
\ai{For a circuit containing $M$ gates, the output is the ordered composition
$\rho_{\mathrm{out}}=(\mathcal E_{g_M}\!\circ\mathcal U_{g_M})\circ\cdots\circ
(\mathcal E_{g_1}\!\circ\mathcal U_{g_1})(\rho_0)$.  Noise therefore accumulates
with the actual transpiled gate count rather than being applied once to the final
state.  Qiskit Aer implements this construction with a density-matrix simulator
and a \texttt{NoiseModel}~\cite{qiskit2024}.}

\ai{The controlled baseline attaches a two-qubit depolarizing channel after
every \texttt{cx},}
\begin{equation}
  \mathcal E^{\mathrm{dep}}_p(\rho)
  =(1-p)\rho+p\,\frac{\operatorname{Tr}\rho}{2^n}I .
  \label{eq:depol}
\end{equation}
\ai{A second model uses gate-duration-dependent thermal relaxation.  For a gate
of duration $\tau$, amplitude loss has probability
$p_{T_1}=1-e^{-\tau/T_1}$, while $T_2$ controls the decay of off-diagonal
coherences.  The two mechanisms differ operationally: relaxation may move the
prepared state out of its target parity sector, whereas dephasing directly
suppresses the endpoint coherence needed by the edge string.  These are channel
mechanics only; their classification as physical errors of the Kitaev chain is
deferred to Part 6.}

\subsection{Three locations for error}

\ai{The study separates effects that can otherwise be confused.  The Week-8
\emph{frozen-state} test takes the validated state and applies readout or a single
post-preparation channel.  Symmetric assignment error $\epsilon$ attenuates a
weight-$L$ Pauli string as}
\begin{equation}
  \langle O_{\mathrm{edge}}\rangle_{\mathrm{readout}}
  =(1-2\epsilon)^L\langle O_{\mathrm{edge}}\rangle,
  \label{eq:readout}
\end{equation}
\ai{with an additional statistical uncertainty of order
$N_{\mathrm{shots}}^{-1/2}$.  The \emph{circuit-level} test instead applies a
channel after every entangling gate.  Finally, a coherent calibration error
replaces $\theta^*$ by $\theta^*+\delta\theta$ but keeps the state pure; an
incoherent channel generally lowers $\operatorname{Tr}\rho^2$.  Purity therefore
distinguishes a unitary miscalibration from genuine mixing even when both reduce
$|\langle O_{\mathrm{edge}}\rangle|$.}

\subsection{Expressibility versus accumulated noise}

\ai{For the linear-entanglement \texttt{EfficientSU2} ansatz, $r$ repetitions on
$L$ qubits contain $r(L-1)$ CNOTs.  A useful approximation to the measured
signal is}
\begin{equation}
 |\langle O_{\mathrm{edge}}\rangle|_{\mathrm{meas}}(r)
 \simeq
 |\langle O_{\mathrm{edge}}\rangle|_{\mathrm{ideal}}(r)
 (1-p_{\mathrm{cx}})^{r(L-1)} .
 \label{eq:depth_tradeoff}
\end{equation}
\ai{The first factor can have an expressibility threshold; the second decreases
at every added layer.  The relevant optimum is consequently the shallowest
depth that can represent the target state, not the deepest affordable circuit.}

\ai{Putting noise inside the VQE changes the objective itself.  With parity
$P=\prod_j Z_j$, the cost evaluated on the noisy density matrix is}
\begin{equation}
 C_{\mathrm{noisy}}(\theta)
 =\operatorname{Tr}[H\rho_\theta]
 +\lambda\bigl(1-\operatorname{Tr}[P\rho_\theta]\bigr)^2,
 \qquad
 \rho_\theta=\mathcal E\!\left(U(\theta)|0\rangle\langle0|U^\dagger(\theta)\right).
 \label{eq:noisy_cost}
\end{equation}
\ai{A lower value of Eq.~\eqref{eq:noisy_cost}, or even a larger noisy edge
string, need not imply a higher-fidelity ground state.  The re-optimized angles
must therefore be checked in a noiseless circuit against the exact ground state.}
